\chapter{低温工业等离子体物理}
\label{ch:low-temp-industrial-plasma}

\begin{chapterintro}
低温等离子体（$T_e \sim 1\sim 10\,\mathrm{eV}$，$T_i \sim$ 室温）虽然温度远低于核聚变等离子体，却是现代半导体工业、材料加工、表面处理的核心工具。本章面向自学者，从气体放电的物理基础出发，深入讲解鞘层这一被全书此前章节忽略的关键结构，再延伸到化学反应、磁控溅射与工艺参数估算。我们的目标是：不仅知道公式，更能在脑中"看到"等离子体在反应腔中的运动图像。
\end{chapterintro}

\section{气体电离与放电基础}
\difficulty{基础}

\subsection{气体电离与汤森放电机制}

当一个中性气体被置于足够强的电场中时，少数天然存在的"种子电子"（例如来自宇宙射线）被加速，获得足以电离中性原子的动能。电离产生的新电子又继续加速，形成雪崩式的链式反应——这便是汤森放电（Townsend discharge）的物理图像。

\begin{definition}{汤森第一电离系数}{}{}
\label{def:townsend-alpha}
汤森第一电离系数 $\alpha$ 定义为电子在电场方向每前进单位距离所产生的平均电离碰撞次数，单位 $\mathrm{m^{-1}}$。它满足经验公式
\begin{equation}
\label{eq:townsend-alpha}
\frac{\alpha}{p} = A \,\exp\!\left(-\frac{Bp}{E}\right)
\end{equation}
其中 $p$ 为气压，$E$ 为电场强度，$A$、$B$ 为与气体种类相关的实验常数。
\end{definition}

\begin{insight}{}{}
\cref{def:townsend-alpha} 的物理含义非常直观：气压 $p$ 越高，电子自由程越短，所以每次电离需要更强的电场 $E$ 来补偿（指数中 $Bp/E$ 项）；但一旦电场足够强，电离系数随 $E/p$ 呈指数增长，雪崩便不可阻挡。
\end{insight}

\begin{history}{}{}
1903年，英国物理学家约翰·汤森（John Townsend）系统研究了低气压气体中的电流增长规律，建立了电子雪崩理论。他用手摇真空泵和简单电极做出了精密的电离系数测量，奠定了气体放电物理的基石。
\end{history}

\begin{theorem}{汤森放电击穿条件}{}{}
\label{thm:townsend-breakdown}
考虑平行板电极间距为 $d$，施加电压 $V$。当满足
\begin{equation}
\label{eq:townsend-breakdown}
\gamma \,\bigl(\ee^{\alpha d} - 1\bigr) = 1
\end{equation}
时，气体发生自持放电。其中 $\gamma$ 为汤森第二电离系数（每个入射离子在阴极产生的二次电子数）。
\end{theorem}

\begin{insight}{}{}
\cref{thm:townsend-breakdown} 的左边 $\gamma(\ee^{\alpha d}-1)$ 表示一个电子从阴极出发、在阳极被收集前，通过电离雪崩在阴极"生出"的新电子数。如果这个数等于1，系统就实现了自持：不需要外部种子电子，放电也能维持。这是反馈放大器中的"正反馈"概念在等离子体中的最早体现。
\end{insight}

\begin{warning}{}{}
初学者常混淆 $\alpha$ 与 $\gamma$ 的物理来源：$\alpha$ 是电子与中性气体碰撞产生的体电离，发生在气体内部；$\gamma$ 是离子轰击阴极产生的表面二次电子发射，发生在电极表面。二者是完全不同的物理过程！
\end{warning}

\subsection{帕邢定律：击穿电压与气压的"U型"关系}

将 \cref{eq:townsend-alpha} 代入 \cref{eq:townsend-breakdown}，可以得到击穿电压 $V_{\text{bd}}$ 与气压 $p$ 和极距 $d$ 的乘积 $pd$ 之间的普遍关系——帕邢定律。

\begin{theorem}{帕邢定律}{}{}
\label{thm:paschen-law}
气体的击穿电压仅取决于 $pd$ 的乘积：
\begin{equation}
\label{eq:paschen-law}
V_{\text{bd}} = \frac{B\,pd}{\ln(Apd) - \ln\!\bigl[\ln(1+1/\gamma)\bigr]}
\end{equation}
对于给定气体和电极材料，$V_{\text{bd}}$ 随 $pd$ 的变化呈U型曲线，存在最小击穿电压 $V_{\text{min}}$。
\end{theorem}

\begin{insight}{}{}
帕邢曲线的U型非常直观：
\begin{itemize}
\item \textbf{左侧（低 $pd$）}：气压太低或极距太小，电子自由程过长，还没撞到中性原子就打到阳极去了，电离效率低，需要很高电压才能击穿。
\item \textbf{右侧（高 $pd$）}：气压太高，电子自由程太短，每次碰撞能量不足以电离，需要更强电场补偿，击穿电压再次升高。
\item \textbf{最低点}：存在最优的自由程，使电子在两次碰撞间刚好积累到电离阈值能量。对于空气在标准电极间距下，$V_{\text{min}} \approx 330\,\mathrm{V}$。
\end{itemize}
\end{insight}

\begin{example}{霓虹灯管中的帕邢曲线应用}{}{}
\label{ex:neon-paschen}
霓虹灯管内充入约 $p = 1\sim 10\,\mathrm{Torr}$ 的氖气，管径 $d \approx 10\,\mathrm{mm}$。计算 $pd = 10\sim 100\,\mathrm{Torr\cdot mm}$，恰好位于帕邢曲线最低点右侧附近。这是霓虹灯只需数百伏即可点亮的根本原因。

\textbf{数值估算}：取氖气 $A = 4\,\mathrm{m^{-1}Pa^{-1}}$，$B = 150\,\mathrm{V\,m^{-1}Pa^{-1}}$，$\gamma \approx 0.05$，$p = 5\,\mathrm{Torr} = 667\,\mathrm{Pa}$，$d = 0.02\,\mathrm{m}$，则 $pd = 13.3\,\mathrm{Pa\,m}$。代入 \cref{eq:paschen-law} 估算得 $V_{\text{bd}} \approx 400\sim 500\,\mathrm{V}$，与商用霓虹灯启辉电压一致。
\end{example}

\subsection{DC辉光放电：从阴极到阳极的"分区宇宙"}

当汤森放电过渡到辉光放电（glow discharge），电流达到 $\sim 1\,\mathrm{mA}$ 量级，气体开始显著发光。平行板间的空间呈现出明显的空间结构，每个区域都有独特的物理特征。

\begin{definition}{DC辉光放电的空间结构}{}{}
\label{def:glow-regions}
从阴极到阳极，典型的辉光放电可分为：
\begin{enumerate}
\item \textbf{阴极暗区（Cathode dark space）}：电子刚从阴极二次发射，能量低，不足以激发原子发光，空间呈现暗区。电场最强，主要压降集中于此。
\item \textbf{负辉光（Negative glow）}：电子能量达到激发阈值，产生强烈发光。这是辉光放电中最亮的区域，等离子体密度最高。
\item \textbf{法拉第暗区（Faraday dark space）}：电子能量因多次非弹性碰撞降低，再次低于激发阈值，发光减弱。
\item \textbf{正柱区（Positive column）}：电场再次加速电子，能量回升，形成准中性等离子体柱，是辉光放电的主体。颜色取决于气体种类（氖气发红光，氩气发紫蓝光）。
\item \textbf{阳极辉光（Anode glow）}：阳极附近电子密度梯度引起的微弱发光。
\end{enumerate}
\end{definition}

\begin{insight}{}{}
可以把DC辉光放电想象成一个"电子过山车"：阴极暗区是加速爬坡（强电场），负辉光是高速俯冲时尖叫的顶点（能量刚好达到激发阈值），法拉第暗区是能量耗尽的平缓段，正柱区是再次加速后的匀速行驶。整个过程中电子在电场中"冲浪"，不断经历加速、碰撞、激发、再加速的循环。
\end{insight}

\begin{warning}{}{}
常见误区：认为负辉光最亮所以电场最强。实际上，负辉光区电场已经较弱，亮度来源于大量电子能量恰好落在激发函数峰值附近；而阴极暗区电场最强，反而不发光，因为电子能量还未达到激发阈值。亮度和电场强度没有直接对应关系！
\end{warning}

\subsection{射频放电：工业应用的主力军}

DC辉光放电有一个致命缺陷：绝缘靶材或介质窗上会积累电荷，迅速熄灭放电。工业上广泛采用射频（RF）放电来克服这一问题。

\begin{definition}{射频放电的频率窗口}{}{}
\label{def:rf-frequency}
工业射频等离子体主要使用 $13.56\,\mathrm{MHz}$（ISM频段，由国际电信联盟分配）。该频率满足两个条件：
\begin{enumerate}
\item 远高于离子等离子体频率（$f_{pi} \sim 1\,\mathrm{MHz}$），离子来不及响应射频电场的变化；
\item 远低于电子等离子体频率（$f_{pe} \sim 1\,\mathrm{GHz}$），电子可以充分响应。
\end{enumerate}
因此，离子看到的是时间平均电场，电子则能形成准稳态的射频鞘层振荡。
\end{definition}

\begin{insight}{}{}
\cref{def:rf-frequency} 的物理图像：在 $13.56\,\mathrm{MHz}$ 的一个周期（$74\,\mathrm{ns}$）内，电子可以跑完德拜长度量级的距离（见第2章），而离子仅移动约 $10^{-4}$ 倍德拜长度。因此鞘层边界以射频频率往复振荡，离子像"慢动作"一样穿过这个振荡边界，被持续加速。这正是射频鞘层可以产生高能离子轰击的关键。
\end{insight}

射频放电有两种主要耦合方式：

\begin{definition}{容性耦合与感性耦合}{}{}
\label{def:ccp-icp}
\begin{itemize}
\item \textbf{容性耦合等离子体（CCP, Capacitively Coupled Plasma）}：射频功率通过平行板电容耦合进入等离子体，电极既是电容极板也是反应边界。电场沿垂直于电极方向，适合薄膜沉积和均匀刻蚀。
\item \textbf{感性耦合等离子体（ICP, Inductively Coupled Plasma）}：射频功率通过线圈电感耦合进入，利用交变磁场感应出涡旋电场维持放电。功率沉积发生在体等离子体区域，等离子体密度可达 $10^{18}\sim 10^{19}\,\mathrm{m^{-3}}$，远高于CCP（$10^{15}\sim 10^{17}\,\mathrm{m^{-3}}$）。
\end{itemize}
\end{definition}

\subsection{微波放电：电子回旋共振（ECR）}

当频率进入微波波段（$2.45\,\mathrm{GHz}$），电磁波的波长与腔体尺寸相当，可以形成驻波模式。若同时施加磁场，使电子回旋频率 $\omega_{ce} = eB/m_e$ 与微波频率共振，则发生高效的电子回旋共振（ECR）加热。

\begin{theorem}{ECR共振条件}{}{}
\label{thm:ecr-condition}
电子回旋共振发生的条件是
\begin{equation}
\label{eq:ecr-condition}
\omega_{\text{MW}} = \omega_{ce} = \frac{eB}{m_e}
\end{equation}
对于 $2.45\,\mathrm{GHz}$ 的微波，共振磁场为 $B \approx 875\,\mathrm{G}$（$0.0875\,\mathrm{T}$）。
\end{theorem}

\begin{insight}{}{}
ECR可以产生极高密度（$n_e \sim 10^{18}\sim 10^{19}\,\mathrm{m^{-3}}$）和低离子温度（$T_i \sim 0.1\,\mathrm{eV}$）的等离子体，因为加热只针对电子，离子基本不直接吸收微波能量。ECR源在离子注入、深度反应离子刻蚀（DRIE）等需要高密度、低损伤的工艺中不可替代。
\end{insight}

\begin{miniquiz}
\begin{enumerate}
\item 为什么霓虹灯管的设计要保证 $pd$ 处于帕邢曲线右侧而非左侧？
\item 射频放电为什么使用 $13.56\,\mathrm{MHz}$ 而非任意频率？
\item 比较CCP和ICP：哪种耦合方式更适合高纵横比刻蚀孔洞（需要高离子密度和各向异性）？
\end{enumerate}
\end{miniquiz}

\section{等离子体鞘层详细理论}
\difficulty{提高}

\subsection{鞘层形成的物理图像}

第2章讨论了德拜屏蔽（Debye shielding），即等离子体对局部电扰动的响应长度尺度为德拜长度 $\lambda_D$。然而，全书此前并未深入讨论这一屏蔽效应在固体边界处的最终表现——等离子体鞘层（sheath）。

\begin{definition}{等离子体鞘层}{}{}
\label{def:sheath}
等离子体鞘层是紧邻固体壁面、准中性被破坏（$n_e \neq n_i$）的薄层区域。其典型厚度为 $3\sim 5\,\lambda_D$（见第2章德拜长度定义）。鞘层内存在强电场，将电子反射回等离子体，同时加速离子向壁面运动。
\end{definition}

\begin{insight}{}{}
鞘层可以形象地看作一个"离子加速器"+"电子屏障"的二合一结构：
\begin{itemize}
\item \textbf{电子屏障}：鞘层电场方向从等离子体指向壁面，电子（带负电）受到的电场力指向等离子体内部，因此大多数电子被反射回去，只有能量最高的"尾巴"电子能克服势垒到达壁面。
\item \textbf{离子加速器}：离子（带正电）被鞘层电场加速向壁面，获得定向动能。在工艺应用中，这些离子动能决定了溅射产额和刻蚀各向异性。
\end{itemize}
\end{insight}

\begin{theorem}{壁面浮置电位}{}{}
\label{thm:floating-potential}
若壁面电绝缘且没有外部电流回路，则壁面自动达到浮置电位（floating potential）$V_f$，使得到达壁面的电子流与离子流相等：
\begin{equation}
\label{eq:floating-potential}
\frac{1}{4} n_e \bar{v}_e \,\exp\!\left(\frac{eV_f}{k_B T_e}\right) = n_i u_B
\end{equation}
由此解得
\begin{equation}
V_f = -\frac{k_B T_e}{2e} \ln\!\left(\frac{m_i}{2\pi m_e}\right)
\end{equation}
对于氩离子，$V_f \approx -4.7\,\dfrac{k_B T_e}{e}$；即壁面电位比等离子体电位低约 $4\sim 5\,T_e$（以eV为单位）。
\end{theorem}

\begin{warning}{}{}
高频误区：认为壁面电位就是等离子体电位。实际上，由于电子热速度远大于离子声速，壁面必须足够负才能排斥过剩的电子，否则电子流会远大于离子流。任何没有外部电流的绝缘表面都会自动"负漂移"到浮置电位！
\end{warning}

\subsection{玻姆判据：离子进入鞘层的最低速度}

\begin{theorem}{玻姆判据（Bohm Criterion）}{}{}
\label{thm:bohm-criterion}
离子要进入鞘层区域，其在鞘层边界的速度必须至少等于离子声速（Bohm速度）：
\begin{equation}
\label{eq:bohm-criterion}
u_B \ge u_B = \sqrt{\frac{k_B T_e}{m_i}}
\end{equation}
\end{theorem}

\begin{proof}[推导概要（挑战）]
设鞘层边界为 $x=0$，鞘层内部 $x>0$。假设冷离子（$T_i = 0$）和 Boltzmann 分布电子：
\begin{align}
n_i(x) &= n_0 \left(1 - \frac{2e\phi(x)}{m_i u_0^2}\right)^{-1/2} \\
n_e(x) &= n_0 \exp\!\left(\frac{e\phi(x)}{k_B T_e}\right)
\end{align}
其中 $u_0$ 为鞘层边界处离子速度，$n_0$ 为边界密度。准中性要求鞘层边界处 $\diff^2\phi/\diff x^2 \to 0$。对 Poisson 方程
\begin{equation}
\eps_0 \frac{\diff^2\phi}{\diff x^2} = e(n_e - n_i)
\end{equation}
在鞘层边界展开，要求右边对 $\phi$ 的导数 $\ge 0$（否则鞘层无法稳定存在），得到
\begin{equation}
\frac{1}{k_B T_e} \ge \frac{m_i u_0^2}{k_B T_e} \quad\Rightarrow\quad u_0 \ge \sqrt{\frac{k_B T_e}{m_i}} = u_B
\end{equation}
\end{proof}

\begin{insight}{}{}
玻姆判据的物理含义极其深刻：鞘层边界的离子速度不能任意小，否则任何微小的电位扰动都会被等离子体"抹平"（德拜屏蔽太强），无法形成稳定的鞘层。只有当离子具有足够的惯性（速度 $\ge u_B$），才能"冲过"准中性到非准中性的过渡，形成鞘层所需的电荷分离。可以把离子比作"运动员"：只有速度足够快，才能冲出起跑线（预鞘区）进入真正的赛道（鞘层）。
\end{insight}

\subsection{预鞘（Pre-sheath）与准中性过渡}

离子如何获得玻姆速度？答案是：在真正的鞘层之前，存在一个更宽的"预鞘"（pre-sheath）区域。

\begin{definition}{预鞘}{}{}
\label{def:pre-sheath}
预鞘是鞘层边界与准中性体等离子体之间的过渡区域，尺度远大于鞘层厚度（通常为装置尺寸量级）。在预鞘中，等离子体仍近似准中性（$n_e \approx n_i$），但存在微弱的电场使离子从热速度加速到 $u_B$。预鞘中电位下降约为 $\frac{1}{2}k_B T_e/e$。
\end{definition}

\begin{insight}{}{}
预鞘的物理图像是：体等离子体内部没有净电场，但离子以热速度随机运动；当离子接近鞘层边界时，感受到来自鞘层"边缘"的微弱"拖曳"，逐渐获得定向速度。这个加速区延伸到约装置尺寸 $L$ 量级，因为准中性要求任何电场都必须足够弱，使得 $n_e$ 和 $n_i$ 的微小差异恰好由空间电荷方程平衡。在预鞘末端，离子终于达到 $u_B$，然后进入真正的鞘层，被急剧加速。
\end{insight}

\begin{history}{}{}
1949年，英国理论物理学家大卫·玻姆（David Bohm）在研究电弧放电时提出了离子进入鞘层的最低速度判据。这一简洁而深刻的结论成为此后70年等离子体鞘层理论的基石，深刻影响了所有低温等离子体工艺的设计。
\end{history}

\subsection{射频鞘层的时变特性}

在射频放电中，鞘层不是静态的，而是随射频电压周期性膨胀和收缩。

\begin{theorem}{射频鞘层时变特性}{}{}
\label{thm:rf-sheath}
对于容性耦合射频放电，鞘层厚度 $s(t)$ 随时间振荡。在一个射频周期内：
\begin{itemize}
\item 当电极瞬时电位最负时，鞘层最厚（最大排斥电子）；
\item 当电极瞬时电位最正时，鞘层最薄（甚至暂时消失，电子短暂涌入）。
\end{itemize}
离子由于惯性大，仅响应时间平均电位 $\bar{V}$，获得的能量为
\begin{equation}
\label{eq:ion-energy-rf}
E_i \approx e\bar{V} + \frac{1}{2}k_B T_i \approx e\bar{V}
\end{equation}
其中 $\bar{V} \sim V_{\text{RF}}/2$，$V_{\text{RF}}$ 为射频电压幅值。
\end{theorem}

\begin{insight}{}{}
射频鞘层的动态图像是：鞘层边界像一个以 $13.56\,\mathrm{MHz}$ 频率往复运动的"活塞"，电子像被活塞来回拍打的"气体分子"，在鞘层边界附近被压缩和膨胀。离子则像缓慢移动的"弹丸"，穿过这个振荡活塞，最终获得由时间平均鞘层压降决定的动能。这种"离子冲浪"机制是RIE（反应离子刻蚀）中离子能量可控的关键。
\end{insight}

\subsection{Child-Langmuir定律：鞘层厚度与电压的关系}

\begin{theorem}{Child-Langmuir定律（无碰撞鞘层）}{}{}
\label{thm:child-langmuir}
对于无碰撞鞘层，假设离子以玻姆速度进入鞘层，电子被完全排斥，则鞘层厚度 $s$ 与鞘层压降 $V_s$ 的关系为
\begin{equation}
\label{eq:child-langmuir}
s = \frac{2}{3}\sqrt{\frac{4\eps_0}{9}} \left(\frac{2e}{m_i}\right)^{1/4} \frac{V_s^{3/4}}{(n_s u_B)^{1/2}}
\end{equation}
化简后
\begin{equation}
s \approx \lambda_D \left(\frac{2V_s}{T_e}\right)^{3/4}
\end{equation}
其中 $V_s$ 以伏特为单位，$T_e$ 以eV为单位。典型地，$V_s \sim 5\,T_e$，则 $s \sim 3\sim 5\,\lambda_D$。
\end{theorem}

\begin{proof}[推导（挑战）]
在无碰撞鞘层中，离子满足能量守恒和运动连续性：
\begin{align}
\frac{1}{2}m_i u^2(x) + e\phi(x) &= \frac{1}{2}m_i u_B^2 \quad (\text{能量守恒}) \\
n_i(x) u(x) &= n_s u_B \quad (\text{连续性})
\end{align}
其中 $n_s$ 为鞘层边界密度。电子被完全排斥，$n_e = 0$。Poisson方程变为
\begin{equation}
\frac{\diff^2\phi}{\diff x^2} = -\frac{e n_i}{\eps_0}
\end{equation}
将 $n_i$ 用 $\phi$ 表示，积分两次，并应用边界条件 $\phi(0)=0$，$\diff\phi/\diff x(0)=0$（准中性过渡），$\phi(s)=-V_s$，最终得到 \cref{eq:child-langmuir}。
\end{proof}

\begin{insight}{}{}
Child-Langmuir定律的 $V_s^{3/4}$ 依赖关系非常有趣：鞘层厚度不是线性随电压增长，而是增长得更慢（因为更强的电压加速离子更快，使离子密度降低，部分抵消了空间电荷效应）。实际工艺中，可以通过调节射频功率来改变鞘层压降，从而控制离子轰击能量，但鞘层厚度始终维持在几个德拜长度——这是德拜屏蔽的普适结果。
\end{insight}

\begin{warning}{}{}
常见误区：认为鞘层厚度可以任意大。实际上，由于离子进入鞘层的速度被玻姆判据固定，且密度随电压升高而下降，鞘层厚度被锁定在德拜长度量级的 $3\sim 5$ 倍。即使功率增加10倍，鞘层厚度可能仅增加不到1倍。这也是射频鞘层可以维持高度各向异性刻蚀的原因——鞘层很薄，离子几乎垂直入射。
\end{warning}

对于高气压（有碰撞）鞘层，离子在鞘层内经历电荷交换碰撞，有效离子质量增加，鞘层变厚。这种情况下需使用修正的 Child-Langmuir 定律或流体模拟，此处从略 [挑战]。

\begin{miniquiz}
\begin{enumerate}
\item 为什么绝缘壁面会自动达到负电位？从电子和离子流平衡角度解释。
\item 玻姆判据中 $u_B = \sqrt{k_B T_e/m_i}$ 的分子为什么是 $T_e$ 而非 $T_i$？
\item 射频鞘层振荡时，为什么离子能量由时间平均电位决定而非瞬时电位？
\end{enumerate}
\end{miniquiz}

\section{等离子体化学反应与工艺应用}

\subsection{活性物种：自由基、离子、激发态}

低温等离子体中的化学反应由三种主要活性物种驱动：

\begin{definition}{活性物种}{}{}
\label{def:active-species}
\begin{itemize}
\item \textbf{自由基（Radicals）}：中性碎片，如 $\mathrm{CF}_x$、$\mathrm{Cl}$、$\mathrm{F}$。寿命长（ms量级），扩散远，主导各向同性化学刻蚀。
\item \textbf{离子（Ions）}：带电粒子，如 $\mathrm{Ar}^+$、$\mathrm{CF}_x^+$。受鞘层电场加速，定向轰击表面，主导各向异性物理效应。
\item \textbf{激发态（Excited states）}：电子跃迁到高能级的原子或分子，如 $\mathrm{Ar}^*$。通过辐射退激发光（辉光），或将能量传递给表面引发化学反应。
\end{itemize}
\end{definition}

\begin{insight}{}{}
三种物种的"分工"非常清晰：自由基像"化学溶剂"，到处扩散，溶解材料（化学刻蚀）；离子像"手术刀"，定向切割，实现精密图形转移（物理+化学刻蚀）；激发态则像"信使"，传递能量但不直接参与反应。三者的协同是等离子体工艺的核心。
\end{insight}

\subsection{刻蚀机理：物理 vs 化学 vs 反应离子刻蚀}

\begin{definition}{三种刻蚀机制}{}{}
\label{def:etching-mechanisms}
\begin{itemize}
\item \textbf{物理溅射刻蚀（Sputtering）}：高能离子（$E_i > 50\,\mathrm{eV}$）直接轰击靶材，通过动量传递将表面原子撞出。纯物理过程，各向异性好，但选择性差（刻蚀速率对材料种类不敏感），且可能损伤下层结构。
\item \textbf{化学刻蚀（Chemical etching）}：自由基与表面发生化学反应，生成挥发性产物（如 $\mathrm{Si} + 4\mathrm{F} \to \mathrm{SiF}_4\uparrow$）。各向同性（自由基各向扩散），选择性好，但无法做精细图形。
\item \textbf{反应离子刻蚀（RIE, Reactive Ion Etching）}：物理溅射与化学刻蚀的结合。离子轰击破坏表面化学键，增强自由基反应；同时反应产物被离子轰击移除。通过调节离子能量和自由基密度，实现各向异性、高选择性的可控刻蚀。
\end{itemize}
\end{definition}

\begin{theorem}{RIE的协同效应}{}{}
\label{thm:rie-synergy}
在RIE中，离子轰击与化学反应的协同使刻蚀速率远大于纯物理和纯化学刻蚀速率之和：
\begin{equation}
\label{eq:rie-synergy}
R_{\text{RIE}} \gg R_{\text{phys}} + R_{\text{chem}}
\end{equation}
这一现象称为离子增强刻蚀（Ion-enhanced etching），其微观机制包括：离子轰击破坏表面钝化层、提高表面温度、增加反应位点。
\end{theorem}

\begin{insight}{}{}
RIE的协同效应是等离子体工艺中最精妙的"化学反应工程"：想象一个顽固的油漆表面，纯化学溶剂（自由基）渗透慢，纯砂纸（离子溅射）效率低且伤底材；但如果一边用溶剂软化，一边用砂纸打磨，效率就能提高数倍。这正是RIE中离子"催化"化学反应的物理图像。
\end{insight}

\begin{example}{手机芯片制造中的等离子体刻蚀}{}{}
\label{ex:chip-etching}
现代手机芯片（7 nm、5 nm工艺）中，晶体管栅极宽度仅几十纳米，需要刻蚀深宽比 $> 50$ 的硅通孔（TSV）。工艺过程如下：
\begin{enumerate}
\item 在硅片上旋涂光刻胶，通过极紫外光刻（EUV）曝光形成图形；
\item 放入ICP-RIE刻蚀机：$\mathrm{SF}_6$ 产生 $\mathrm{F}$ 自由基，$\mathrm{C}_4\mathrm{F}_8$ 产生聚合物钝化层；
\item 射频鞘层加速 $\mathrm{SF}_x^+$ 离子垂直轰击底部，移除底部聚合物，同时 $\mathrm{F}$ 自由基化学刻蚀暴露的硅；
\item 侧壁因无离子轰击而被聚合物保护，实现高度各向异性；
\item 通过 "Bosch工艺" 循环切换刻蚀与沉积步骤，实现高深宽比刻蚀。
\end{enumerate}
\textbf{数值估算}：ICP功率 $2\,\mathrm{kW}$，$n_e \sim 10^{17}\,\mathrm{m^{-3}}$，鞘层压降 $200\,\mathrm{V}$，离子能量 $E_i \sim 200\,\mathrm{eV}$，硅的溅射阈值约 $20\,\mathrm{eV}$，化学刻蚀速率约 $1\,\mu\mathrm{m}/\mathrm{min}$，物理溅射速率约 $0.1\,\mu\mathrm{m}/\mathrm{min}$，RIE协同后速率可达 $5\sim 10\,\mu\mathrm{m}/\mathrm{min}$。
\end{example}

\begin{warning}{}{}
常见误区：认为高深宽比刻蚀只要离子能量足够高就行。实际上，侧壁保护和钝化层管理才是关键。纯物理溅射在深孔底部会产生"微沟槽效应"（micro-trenching），因为离子在侧壁反射后汇聚到底部角落；纯化学刻蚀则完全各向同性。Bosch工艺的成功在于精确平衡了刻蚀与钝化的循环，而非简单提高功率。
\end{warning}

\subsection{薄膜沉积：PECVD与溅射沉积}

\begin{definition}{等离子体增强化学气相沉积（PECVD）}{}{}
\label{def:pecvd}
PECVD利用等离子体中的电子碰撞将反应气体（如 $\mathrm{SiH}_4$、$\mathrm{NH}_3$）分解为活性自由基，在衬底表面发生化学反应形成薄膜。等离子体的作用是将热激活能 $> 500\,\mathrm{^\circ C}$ 的CVD反应降低到 $200\sim 300\,\mathrm{^\circ C}$，避免高温对已有器件的损伤。
\end{definition}

\begin{definition}{溅射沉积}{}{}
\label{def:sputter-deposition}
溅射沉积利用辉光放电或磁控放电产生的离子轰击靶材，将靶原子溅射出来并沉积到衬底上形成薄膜。磁控溅射（见 \cref{sec:magnetron-sputtering}）通过磁场约束提高等离子体密度和沉积速率。
\end{definition}

\begin{insight}{}{}
PECVD与溅射沉积的物理机制完全不同：PECVD是"化学反应在表面"，等离子体只在气相中生成前驱体；溅射沉积是"物理动量传递"，靶材原子被"撞飞"后沉积。因此PECVD薄膜的化学计量比可以灵活调控（改变气体配比），而溅射沉积薄膜的组分与靶材基本一致。实际工艺中，PECVD用于 $\mathrm{SiO}_2$、$\mathrm{SiN}_x$ 等绝缘层，溅射用于金属互连（Al、Cu、W）。
\end{insight}

\begin{miniquiz}
\begin{enumerate}
\item 为什么RIE中离子轰击能增强化学刻蚀速率？列举两种微观机制。
\item 比较PECVD与热CVD的温度优势：PECVD为什么可以在低温下沉积？
\item 在Bosch工艺中，侧壁保护靠什么实现？底部刻蚀靠什么实现？
\end{enumerate}
\end{miniquiz}

\section{磁控溅射与刻蚀机物理模型}
\label{sec:magnetron-sputtering}

\subsection{磁控溅射的磁场设计：E\times B漂移}

磁控溅射是工业薄膜沉积的主流技术，其核心物理是利用磁场约束电子，提高等离子体密度和电离效率。

\begin{theorem}{磁控溅射中的E\times B漂移}{}{}
\label{thm:exb-magnetron}
在磁控溅射靶面附近，存在平行于靶面的磁场 $\vect{B}$ 和垂直于靶面的电场 $\vect{E}$（由鞘层电位降产生）。电子在电磁场中做 $E\times B$ 漂移（见第3章 \cref{sec:exb-drift}）：
\begin{equation}
\label{eq:exb-magnetron}
\vect{v}_{E\times B} = \frac{\vect{E}\times\vect{B}}{B^2}
\end{equation}
漂移方向平行于靶面，形成闭合的环形漂移轨道。电子被"困"在靶面附近，大大增加了电离碰撞次数，使等离子体密度提高 $10\sim 100$ 倍。
\end{theorem}

\begin{insight}{}{}
磁控溅射的物理图像非常优美：电子像被磁场"绳子"绑在靶面上，沿着电场和磁场构成的"跑道"不停地跑圈（$E\times B$ 漂移）。每跑一圈，电子在电场中加速，积累能量，然后与气体碰撞产生电离。没有磁场时，电子很快打到阳极或壁面；有磁场时，电子的寿命延长，电离效率暴增。这就像把蜡烛放在四面镜子围成的房间里，光线被不断反射，亮度远超单支蜡烛。
\end{insight}

\begin{definition}{磁控溅射的磁场构型}{}{}
\label{def:magnetron-config}
常见构型包括：
\begin{itemize}
\item \textbf{平面磁控（Planar magnetron）}：圆柱形永磁体在靶后形成拱形磁场线，靶面为平面。$E\times B$ 漂移形成闭合的跑道形（racetrack）等离子体环，溅射也集中在跑道区域，造成靶材不均匀消耗（"跑道侵蚀"）。
\item \textbf{旋转圆柱磁控（Rotating cylindrical magnetron）}：靶材为可旋转的圆筒，内部放置磁体。旋转使靶材均匀消耗，适合大面积连续沉积（如建筑玻璃镀膜）。
\item \textbf{平衡/非平衡磁控}：平衡磁控中磁场线在靶面闭合，等离子体约束在靶面附近；非平衡磁控中部分磁场线延伸到衬底区域，使等离子体同时到达衬底，提高沉积时的离子辅助效果。
\end{itemize}
\end{definition}

\begin{warning}{}{}
常见误区：认为磁控溅射中磁场是用来加速离子轰击靶材的。实际上磁场主要约束电子，离子因质量大几乎不受磁场影响（回旋半径 $\rho_i \gg$ 装置尺寸）。离子是被鞘层电场加速的！磁场通过约束电子来提高等离子体密度，间接增加离子密度，从而增加溅射速率。混淆磁场对电子和离子的作用是初学者常见错误。
\end{warning}

\subsection{刻蚀腔体中的粒子输运：双极扩散到边界损失}

第8章介绍了双极扩散（ambipolar diffusion）——由于电子扩散快于离子，电子先离开会在局部产生电荷分离电场，该电场反过来加速离子、拖慢电子，使二者以相同的"双极扩散系数" $D_a$ 共同扩散。

\begin{theorem}{刻蚀机中的双极扩散损失}{}{}
\label{thm:ambipolar-etcher}
在刻蚀腔体中，等离子体通过双极扩散向壁面和电极损失，损失通量为
\begin{equation}
\label{eq:ambipolar-loss}
\Gamma_{\text{loss}} = D_a \,|\nabla n| \sim \frac{D_a n}{L_{\text{eff}}}
\end{equation}
其中 $L_{\text{eff}}$ 为有效特征长度（与腔体尺寸和几何形状有关）。双极扩散系数（见第8章）为
\begin{equation}
D_a = \frac{k_B T_e}{e} \mu_i \approx \frac{k_B T_e}{m_i \nu_{in}}
\end{equation}
其中 $\mu_i$ 为离子迁移率，$\nu_{in}$ 为离子-中性碰撞频率。损失通量必须与体电离率平衡，以维持稳态等离子体。
\end{theorem}

\begin{insight}{}{}
双极扩散的物理图像是：电子像"急性的兔子"，想迅速逃离；离子像"慢性的乌龟"，跟不上。但兔子一跑远，就被"电场绳子"拽回来（因为电荷分离产生的电场将兔子拉回）。最终兔子和乌龟以中间速度共同移动——这个速度由乌龟（离子）的 mobility 和兔子的"焦躁程度"（$T_e$）共同决定。在刻蚀腔中，粒子不断从中心产生，向边界扩散损失，形成密度分布 $n(r)$。
\end{insight}

\begin{example}{刻蚀腔体中的粒子输运估算}{}{}
\label{ex:etcher-transport}
考虑一个半径 $R = 0.15\,\mathrm{m}$、高度 $h = 0.05\,\mathrm{m}$ 的圆柱形CCP刻蚀腔，充入 $p = 10\,\mathrm{mTorr}$ 的 $\mathrm{Cl}_2$ 气体。估算双极扩散系数和粒子损失时间。

\textbf{解}：
\begin{enumerate}
\item 离子-中性碰撞频率：$\nu_{in} = n_n \sigma v_{th,i}$，其中 $n_n = p/k_B T_n \approx 3.2\times 10^{20}\,\mathrm{m^{-3}}$（$T_n = 300\,\mathrm{K}$），$\sigma \sim 10^{-18}\,\mathrm{m^2}$，$v_{th,i} = \sqrt{k_B T_i/m_i} \sim 300\,\mathrm{m/s}$。得 $\nu_{in} \sim 10^5\,\mathrm{s^{-1}}$。
\item 双极扩散系数：$D_a = k_B T_e/(m_i \nu_{in}) \sim 1.38\times 10^{-23} \times 3\,\mathrm{eV} / (1.6\times 10^{-19} \times 70\times 10^{-3} \times 10^5) \approx 0.04\,\mathrm{m^2/s}$（取 $T_e = 3\,\mathrm{eV}$，$\mathrm{Cl}_2$ 离子质量 $70\,\mathrm{amu}$）。
\item 径向特征长度 $L_{\text{eff}} \sim R/2.4 \approx 0.06\,\mathrm{m}$（圆柱腔最低径向模）。
\item 扩散损失时间：$\tau \sim L_{\text{eff}}^2/D_a \sim (0.06)^2/0.04 \approx 0.09\,\mathrm{s} = 90\,\mathrm{ms}$。
\end{enumerate}

即粒子在腔内平均停留约 $100\,\mathrm{ms}$ 后扩散到壁面损失。要维持稳态，电离产生的速率必须与此损失平衡。
\end{example}

\begin{history}{}{}
1974年，美国物理学家彼得·克兰（Peter J. Clarke）和约翰·帕滕（John S. Chapin）分别独立发明了平面磁控溅射技术。利用永磁体在靶面产生拱形磁场，使电子做 $E\times B$ 漂移，溅射速率提高了10倍以上。这一发明彻底改变了薄膜工业，使低成本大面积镀膜成为可能，从建筑玻璃到硬盘磁头，无处不在。
\end{history}

\begin{miniquiz}
\begin{enumerate}
\item 磁控溅射中，为什么电子被约束而离子基本不受磁场影响？从回旋半径比较说明。
\item 双极扩散系数 $D_a$ 为什么远大于离子自由扩散系数 $D_i$，但远小于电子自由扩散系数 $D_e$？
\item 非平衡磁控溅射与平衡磁控溅射相比，对衬底薄膜质量有何影响？
\end{enumerate}
\end{miniquiz}

\section{低温等离子体特征参数估算案例}

\subsection{典型参数范围}

\begin{table}[htbp]
\centering
\caption{低温等离子体典型参数对比}
\label{tab:plasma-params}
\begin{tabular}{lccc}
\toprule
\textbf{参数} & \textbf{辉光放电} & \textbf{刻蚀机（CCP/ICP）} & \textbf{托卡马克（聚变）} \\
\midrule
$n_e\;(\mathrm{m^{-3}})$ & $10^{16}\sim 10^{18}$ & $10^{16}\sim 10^{18}$（CCP）& $10^{19}\sim 10^{20}$ \\
& & $10^{18}\sim 10^{19}$（ICP）& \\
$T_e\;(\mathrm{eV})$ & $1\sim 10$ & $3\sim 5$ & $10^3\sim 10^4$ \\
$T_i\;(\mathrm{K})$ & $300\sim 500$ & $300\sim 800$ & $10^8$ \\
$\lambda_D\;(\mathrm{m})$ & $10^{-4}\sim 10^{-3}$ & $10^{-4}\sim 10^{-5}$ & $10^{-5}\sim 10^{-4}$ \\
$\omega_{pe}\;(\mathrm{rad/s})$ & $10^8\sim 10^9$ & $10^8\sim 10^9$ & $10^{12}\sim 10^{13}$ \\
$\nu_{en}\;(\mathrm{s^{-1}})$ & $10^7\sim 10^9$ & $10^6\sim 10^8$ & $10^4\sim 10^5$ \\
\bottomrule
\end{tabular}
\end{table}

\begin{insight}{}{}
\cref{tab:plasma-params}揭示了一个深刻对比：聚变等离子体追求 $T_e = T_i \sim 10\,\mathrm{keV}$ 的完全热平衡，而低温等离子体刻意维持 $T_e \gg T_i$（非热平衡）。这种非平衡正是低温等离子体可以"冷处理"热敏感材料的原因：电子提供电离和化学反应能量，离子保持低温，不损伤衬底。这是低温等离子体区别于热等离子体和聚变等离子体的本质特征。
\end{insight}

\subsection{由功率、气压估算等离子体密度}

\begin{example}{CCP刻蚀机中的密度估算}{}{}
\label{ex:density-estimation}
一个容性耦合刻蚀机，功率 $P = 500\,\mathrm{W}$，频率 $13.56\,\mathrm{MHz}$，电极面积 $A = 0.05\,\mathrm{m^2}$，极距 $d = 0.03\,\mathrm{m}$，充入 $p = 5\,\mathrm{mTorr}$ 的氩气。估算等离子体密度。

\textbf{解}：
\begin{enumerate}
\item \textbf{功率密度}：放电体积 $V = A d = 1.5\times 10^{-3}\,\mathrm{m^3}$，功率密度 $P/V = 3.3\times 10^5\,\mathrm{W/m^3}$。
\item \textbf{电离能}：氩的第一电离能 $E_{\text{iz}} = 15.8\,\mathrm{eV}$。假设射频功率中约 $10\sim 30\%$ 有效用于电离，取效率 $\eta = 0.15$。
\item \textbf{电离速率}：每产生一个电子-离子对需要能量 $E_{\text{iz}}$。有效电离功率 $P_{\text{iz}} = \eta P = 75\,\mathrm{W}$。电离速率
\begin{equation}
R_{\text{iz}} = \frac{P_{\text{iz}}}{e E_{\text{iz}}} = \frac{75}{1.6\times 10^{-19} \times 15.8} \approx 3\times 10^{19}\,\mathrm{s^{-1}}
\end{equation}
\item \textbf{损失平衡}：稳态时电离率等于损失率。损失主要由双极扩散到壁面和电极。有效损失面积 $A_{\text{eff}} \approx 2A = 0.1\,\mathrm{m^2}$（上下电极），损失通量 $\Gamma = R_{\text{iz}}/A_{\text{eff}} = 3\times 10^{20}\,\mathrm{m^{-2}s^{-1}}$。
\item \textbf{由玻姆流估算密度}：壁面处离子流为玻姆流 $\Gamma = n_s u_B$，其中 $u_B = \sqrt{k_B T_e/m_i}$。取 $T_e = 4\,\mathrm{eV}$，则 $u_B = \sqrt{1.6\times 10^{-19}\times 4/(40\times 1.67\times 10^{-27})} \approx 3.1\times 10^3\,\mathrm{m/s}$。
\item \textbf{鞘层边界密度}：$n_s = \Gamma/u_B = 3\times 10^{20}/3.1\times 10^3 \approx 1\times 10^{17}\,\mathrm{m^{-3}}$。
\item \textbf{体密度}：预鞘中密度下降约 $1/2$，因此 $n_e \sim 2 n_s \approx 2\times 10^{17}\,\mathrm{m^{-3}}$。
\end{enumerate}

这与 \cref{tab:plasma-params} 中CCP的典型范围一致。若改用ICP（功率 $2\,\mathrm{kW}$，效率更高），密度可提高 $1\sim 2$ 个量级。
\end{example}

\begin{insight}{}{}
\cref{ex:density-estimation} 展示了工程估算的核心思想：不求精确解，而是通过物理平衡关系（功率输入 $=$ 能量损失）和特征尺度（德拜长度、玻姆速度、双极扩散）快速锁定参数数量级。在工业现场，工程师就是这样"估算"等离子体状态的——不是解完整的流体方程，而是抓住主要矛盾（功率平衡、粒子平衡），用"信封背面"算出可用参数。
\end{insight}

\begin{warning}{}{}
常见误区：在密度估算中忘记预鞘密度降。玻姆判据要求鞘层边界密度 $n_s$ 约为体密度 $n_0$ 的一半（预鞘中电位降 $\sim k_B T_e/2e$ 使离子加速、密度下降）。若直接用体密度计算玻姆流，会高估密度约2倍。记住：玻姆流公式中的 $n$ 是鞘层边界密度，不是体密度！
\end{warning}

\begin{miniquiz}
\begin{enumerate}
\item 为什么低温等离子体中 $T_e \gg T_i$ 是有利的？从工艺损伤角度解释。
\item 估算题：一个ICP源功率 $2\,\mathrm{kW}$，$T_e = 5\,\mathrm{eV}$，氩气，电极面积 $0.03\,\mathrm{m^2}$，假设电离效率 $20\%$。估算等离子体密度（提示：先算电离速率，再除以损失面积和玻姆速度）。
\item 比较辉光放电与ICP的参数：为什么ICP密度更高但 $T_e$ 更低？
\end{enumerate}
\end{miniquiz}

\section{本章要点回顾}

\begin{chapterreview}
\begin{enumerate}
\item \textbf{气体放电基础}：汤森放电描述电子雪崩和自持条件；帕邢定律给出击穿电压与 $pd$ 的U型关系；DC辉光放电有明确空间结构（阴极暗区、负辉光、正柱）；射频放电（$13.56\,\mathrm{MHz}$）通过时变电场克服绝缘表面积电荷问题；ECR在 $\omega = \omega_{ce}$ 时实现高效电子加热。

\item \textbf{鞘层理论}：鞘层是壁面附近的非准中性层（$3\sim 5\,\lambda_D$），是"离子加速器"+"电子屏障"；绝缘壁面自动达到浮置电位（比等离子体电位低 $4\sim 5\,T_e$）；玻姆判据要求离子进入鞘层的速度 $\ge u_B = \sqrt{k_B T_e/m_i}$；预鞘是尺度远大于鞘层的准中性过渡区，将离子从热速度加速到 $u_B$；射频鞘层时变振荡，离子响应时间平均电位；Child-Langmuir定律给出鞘层厚度 $s \propto V_s^{3/4}$，锁定在德拜长度量级。

\item \textbf{化学反应与应用}：自由基主导各向同性化学刻蚀，离子主导各向异性物理效应，激发态传递能量；RIE通过离子增强化学刻蚀实现协同效应；PECVD用等离子体降低沉积温度；Bosch工艺通过刻蚀/钝化循环实现高深宽比刻蚀。

\item \textbf{磁控溅射与输运}：磁控溅射利用 $E\times B$ 漂移约束电子（见第3章），提高电离效率；磁场约束电子而非离子；刻蚀腔中粒子通过双极扩散（见第8章）向边界损失，损失通量与体电离率平衡维持稳态。

\item \textbf{参数估算}：辉光放电 $n_e \sim 10^{16}\sim 10^{18}\,\mathrm{m^{-3}}$，$T_e \sim 1\sim 10\,\mathrm{eV}$；CCP刻蚀机 $n_e \sim 10^{16}\,\mathrm{m^{-3}}$，ICP $n_e \sim 10^{18}\,\mathrm{m^{-3}}$；通过功率平衡和玻姆流可以"信封背面"估算等离子体密度。
\end{enumerate}
\end{chapterreview}

\newpage
